\scanpage{22}
% L1–5
Soit $n$ un entier, $k$ un anneau \inserted{comm} avec $2$ inv. ds $k$
$C_n$ catégorie formée des
\[
\left\{\begin{aligned}
& k\text{-module }E\\
& \inserted{1}\ \deleted{\unsure{\text{famille}}}\,\text{famille }(E_i)_{i\in I}\ \text{de sous-modules libres de }E\\
& \forall i\in I,\ \text{partie } B_i\text{ de }E_i
\end{aligned}\right.
\]

% L6–9
avec axiomes
\[
\begin{aligned}
&\text{a) } E\text{ somme directe des } E_i\\
&\text{b) }\operatorname{card} B_i=2\qquad(\forall i\in I)\\
&\text{c) }\forall b\in B_i,\ \{b\}\text{ base de } E_i\\
&\text{d) }B_i=-B_i
\end{aligned}
\]

% L10–15
% slanted
\margin{épure \deleted{\unsure{combinatoire}} \deleted{\unsure{cubique}} de $n$-cube \deleted{\unsure{combinatoire}}}
$\mathcal D_n^{\unsure{*|+}}$ catégorie formée des ens $B$ de card. $2n$,
munis d'une relation d'équiv. $R$ à orbites
de cardinal $2$ (donc $B/R=I$ de card $n$)
-- ou encore des ens. de card $2n$, avec
\deleted{involution} \inserted{$\{\pm1\}$} opérant dessus librement i.e. avec une
involution sans pts fixes.

% L16–20
\[
 C_n\unsure{\longrightarrow|-}\mathcal D_n\qquad\text{par}\quad
 (E,(E_i)_{i\in I},(B_i)_{i\in I})\longmapsto\{\bigcup B_i,\ldots\}
\]
Mais on a aussi
\[
 \mathcal D_n\longrightarrow C_n\qquad\text{en posant}
\]
% The sequence at the left is circled.
\margin{\tikz[baseline=(q.base)]\node[draw,ellipse,inner sep=2pt](q){$\scriptstyle 0\longrightarrow E-k^B\longrightarrow k^I-0$};}
\[
 E=\coprod_{i\in I}B_i\wedge_{\{\pm1\}}k
\]
\margin{\uline{ou aussi} $E\subset k^{(B)}$ formé des $\sum_{b\in B}\lambda_b e_b$ avec $\lambda_{-b}=-\lambda_b$}
où $\{\pm1\}$ opère sur $k$ par homothéties. Donc les
deux catégories sont équivalentes.
% archive sheet number: 21
